% Appendix B (writer: setting). Proofs for sec:versions. Source: research/tracks/foundations/notes-final.md §4 (and notes.md Prop 4.4 for the hashing proof); context Prop 5.4.
\section{Kolmogorov and semimeasure versions: proofs}\label{app:versions}

Notation is that of \cref{sec:versions}. $j(x)$ is the length-lexicographic index of $x\in\{0,1\}^*$, so $j(x)\le2^{d+1}-2$ for $|x|\le d$.

\subsection{Programs that cannot separate inputs}\label{app:versions:Wgap}

\begin{proof}[Proof of \cref{prop:versions:Wgap}]
$\Omtail{\ell}\to0$ because $\sum_{p\in\dom V}2^{-|p|}\le1$.

(a) Let $E_\ell$ be the set of indices $V(p)$ with $|p|<\ell$; it has fewer than $2^\ell$ elements. Colour $x\in\{0,1\}^d$ by $(\chi_e(x))_{e\in E_\ell}$, where $\chi_e(x)=\varphi_e(x)$ if this is $0$ or $1$ and $\chi_e(x)=*$ otherwise. There are at most $3^{2^\ell}$ colours, so if $2^d>3^{2^\ell}$ two distinct strings $x_1,x_2$ share a colour. No $e\in E_\ell$ has $\varphi_e(x_1)=0$ and $\varphi_e(x_2)=1$, so $W(\{(x_1,0),(x_2,1)\})\le\Omtail{\ell}$. Choose $\ell$ with $\Omtail{\ell}\le\varepsilon$, then $d$ with $2^d>3^{2^\ell}$. For the lower bounds, the uniform product law $P(0\mid x)=P(1\mid x)=\frac12$ has complexity $O(1)$ and gives every two-point binary data set probability $\frac14$; so $\Mprod(D)\ge2^{-O(1)}/4$, and $\MF\ge2^{-O(1)}\Mprod$ by \cref{thm:versions:chain}.

(b) Fewer than $2^{d-1}$ programs have length at most $d-2$, and each gives at most one value at $0^d$. So at least $2^{d-1}$ strings $y\in\{0,1\}^d$ receive $W$-mass only from programs of length at least $d-1$. Each program contributes to at most one $y$, so these $W(\{(0^d,y)\})$ sum to at most $\Omtail{d-1}$, and one of them is at most $\Omtail{d-1}/2^{d-1}$. For the lower bound, $P(y\mid x):=2^{-|x|}\ind[|y|=|x|]$ is a computable conditional probability of complexity $O(1)$, so $\Mprod(\{(0^d,y)\})\ge2^{-O(1)}2^{-d}$.
\end{proof}

\begin{proof}[Proof of the claims in \cref{rem:versions:Wrate}]
$\ell_d$ is the largest $\ell$ with $2^\ell\log3<d$, i.e.\ with $3^{2^\ell}<2^d$ (c8(a)). \emph{$\log(1/\Omtail{\ell})\le\Kpad(\ell)+O(1)$}: choose $L\ge\ell$ with $K(L)=\Kpad(\ell)$. The prefix machine that reads a $V$-program for $L$, then $L$ more bits, and halts, is simulated by $V$ with a fixed prefix $r$. The resulting elements of $\dom V$ have length at least $L\ge\ell$ and total weight $2^{L}\cdot2^{-|r|-K(L)-L}$, so $\Omtail{\ell}\ge2^{-\Kpad(\ell)-|r|}$. \emph{Rate}: apply \cref{prop:versions:tape}(c) with $g(N):=h(\max(0,\lfloor\log N\rfloor-1))$, which is computable, nondecreasing and unbounded; then $\Kpad(2^{d+1})\le2\log(h(d)+2)+O_h(1)$.
\end{proof}

\subsection{The tape-length bound and the hashing bound}\label{app:versions:tape}

\begin{proof}[Proof of \cref{prop:versions:tape}]
(a) \emph{Monotonicity and limit.} A prefix of $\omega_{1:L}$ is a prefix of $\omega_{1:L'}$ for $L'\ge L$. The event $f_\omega\supseteq D$ holds exactly when each of the finitely many pairs of $D$ has a certifying triple along $\omega$, and then all of them lie within the longest one; so $\{f_\omega\supseteq D\}=\bigcup_L\{\text{certified within }L\text{ bits}\}$, an increasing union.

\emph{The bound.} Fix $L'\ge L$ with $K(L')=\Kpad(L)$. For $w\in\{0,1\}^{L'}$ let $\varphi_{T,w}(x):=y$ if some $(x,w',y)\in T$ has $w'$ a prefix of $w$, and undefined otherwise. This is well defined because $T$ is consistent and any two prefixes of $w$ are comparable, and it is partial computable (enumerate $T$). Let $V'$ be the prefix machine that reads a $V$-program for an index $i$ of $T$, then a $V$-program for $L'$, then $L'$ further bits $w$, and outputs an index of $\varphi_{T_i,w}$; its domain is prefix-free. By universality $V(rs)=V'(s)$ for a fixed string $r$. With shortest programs $q_i,q_{L'}$, the strings $r\,q_i\,q_{L'}\,w$, $w\in\{0,1\}^{L'}$, are distinct elements of $\dom V$, so
\[
W(D)\ \ge\ \sum_{w\in\{0,1\}^{L'}}2^{-|r|-K(T)-K(L')-L'}\,\ind[\varphi_{T,w}\supseteq D]\ =\ 2^{-|r|-K(T)-\Kpad(L)}\,\nu^{(L')}(D)\ \ge\ 2^{-|r|-K(T)-\Kpad(L)}\,\nu^{(L)}(D).
\]
The consequence for $L_{1/2}$ is immediate, and $K(T_U)=O(1)$.

(b) Apply (a) to the machine $F(x):=\omega_{j(x)+1}$, of complexity $O(1)$, whose law gives every binary-labelled $D$ probability $2^{-|D|}$. For $|x|\le d$ its certificates lie within the first $2^{d+1}-1$ tape bits, and $\Kpad(2^{d+1}-1)\le\Kpad(2^{d+1})$ since $\Kpad$ is nondecreasing. For the second inequality put $s:=\lceil\log(d+1)\rceil$ and $L'':=2^{2^s}\ge2^{d+1}$; then $\Kpad(2^{d+1})\le K(L'')\le K(s)+O(1)$.

(c) $\Kpad$ is nondecreasing because the minimum runs over a shrinking set, and unbounded because only finitely many $L$ have $K(L)\le c$ for each $c$. $\Kpad(2^{2^s})\le K(2^{2^s})\le K(s)+O(1)$. For $g$: let $N^*:=\min\{N':g(N')\ge g(N)+1\}$. It exists because $g$ is unbounded, it exceeds $N$ because $g$ is nondecreasing, and it is computable from $g(N)+1$. So $\Kpad(N)\le K(N^*)\le K(g(N)+1)+c_g\le2\log(g(N)+2)+c'_g$, using $K(m)\le2\log(m+1)+O(1)$.
\end{proof}

\begin{remark}[The hashing bound, superseded]\status{proved; computed}\src{foundations Rem 4.4a; checks c2, r10}\label{rem:versions:hashing}
For every $D$ with labels in $\{0,1\}$, $|D|\le k$ and $\dom D\subseteq\{0,1\}^{\le d}$: $W(D)\ge2^{-|D|-K(d,k)-c}$. \Cref{prop:versions:tape}(b) is stronger, since $K(d,k)\ge K(d)-O(1)\ge\Kpad(2^{d+1})-O(1)$, and it does not depend on $|D|$. The hint in this project's research brief, ``a log factor via $k$-wise independent hashing'', describes the single shortest program, not $W$: for $\KF$ a logarithmic cost is real (\cref{thm:versions:nocoding}), and \cref{prop:versions:covering} bounds it by a covering argument rather than by hashing.
\end{remark}

\begin{proof}
We may assume $k<2^{d+1}$, since $|\{0,1\}^{\le d}|<2^{d+1}$ (replacing $k$ by $\min(k,2^{d+1}-1)$ changes $K(d,k)$ by $O(1)$). Put $q:=d+1$ and fix, computably from $q$, an irreducible polynomial of degree $q$ over $\mathrm{GF}(2)$. Map $x\in\{0,1\}^{\le d}$ injectively to the element $z_x$ of $\mathrm{GF}(2^q)$ whose coefficient vector is the binary string $1x$. For a seed $s\in\mathrm{GF}(2^q)^k$ let $h_s(z)$ be the lowest bit of $\sum_{i<k}s_iz^i$. For distinct $z_1,\dots,z_k$ the Vandermonde map $s\mapsto(\sum_is_iz_l^i)_{l\le k}$ is a bijection of $\mathrm{GF}(2^q)^k$, so for uniform $s$ the values at any $j\le k$ distinct points are uniform and independent, and so are their lowest bits (c2 checks this exhaustively for $(q,k)\in\{(3,2),(3,3),(4,2),(4,3)\}$). Hence $\Prob_s[h_s(z_x)=y\ \forall(x,y)\in D]=2^{-|D|}$. Let $V'$ read a $V$-program for $(d,k)$, then $kq$ further bits $s$, and output an index of $x\mapsto h_s(z_x)$ on $\{0,1\}^{\le d}$; its domain is prefix-free. With $V(rs')=V'(s')$, summing over the $2^{kq}$ seeds gives $W(D)\ge2^{-|r|-K(d,k)}\,\Prob_s[\cdots]=2^{-|r|-K(d,k)-|D|}$.
\end{proof}

\subsection{Shared randomness}\label{app:versions:shared}

\begin{proof}[Proof of \cref{prop:versions:shared}]
$\nulen$ is the law of the function machine $\{(x,w,b):|w|=|x|+1,\ w_{|x|+1}=b\}$, of complexity $O(1)$. Suppose $\Mprod\ge c\,\nulen$ for some $c\in(0,1]$. Put $\varepsilon:=c/8$, write $\Mprod=\sum_iw_i\nu_{P_i}$ with $w_i=\mm(i)$, and choose $J$ with $\sum_{i>J}w_i\le\varepsilon$. For $n\ge1$ let $a_n:=0^n$, $b_n:=1^n$, and $s_i(n,y):=P_i(y\mid a_n)P_i(y\mid b_n)$ for $y\in\{0,1\}$.

\emph{Compactness.} Choose an infinite set of lengths along which $P_i(y\mid a_n)\to\alpha_i(y)$ and $P_i(y\mid b_n)\to\beta_i(y)$ for all $i\le J$ and $y\in\{0,1\}$. Let $S_i:=\alpha_i(0)\beta_i(0)+\alpha_i(1)\beta_i(1)\le\max_y\alpha_i(y)\cdot(\beta_i(0)+\beta_i(1))\le1$. If $S_i=1$, then $\max_y\alpha_i(y)=1$, so $\alpha_i=\delta_{y_i}$ for some $y_i$, and then $\beta_i(y_i)=1$, so $\beta_i=\delta_{y_i}$. Let $I_1:=\{i\le J:S_i=1\}$, $I_0:=\{i\le J:S_i<1\}$ and $r:=\max_{i\in I_0}(1+S_i)/2<1$.

\emph{Subsequence.} Choose lengths $n_1<n_2<\cdots$ from that set such that, for every $j$, $s_i(n_j,0)+s_i(n_j,1)\le r$ for $i\in I_0$, and $s_i(n_j,1-y_i)\le\delta2^{-j}$ for $i\in I_1$, where $\delta:=\varepsilon/2$.

\emph{Data.} For $b\in\{0,1\}^k$ let $D_b:=\{(a_{n_j},b_j),(b_{n_j},b_j):j\le k\}$. Then $\nulen(D_b)=2^{-k}$. All $D_b$ have the same domain, so $\sum_b\nu_{P_i}(D_b)=\prod_j\big(s_i(n_j,0)+s_i(n_j,1)\big)\le1$ for every $i$.

\emph{Count.} Let $B_k$ be $\{0,1\}^k$ without the two constant vectors. For $i\in I_1$ and $b\in B_k$ the set $T_b:=\{j:b_j\ne y_i\}$ is nonempty, so $\prod_js_i(n_j,b_j)\le\prod_{j\in T_b}\delta2^{-j}$, and $b\mapsto T_b$ is injective; hence $\sum_{b\in B_k}\prod_js_i(n_j,b_j)\le\prod_j(1+\delta2^{-j})-1\le e^\delta-1\le2\delta$. Therefore
\[
\textstyle\sum_{b\in B_k}\Mprod(D_b)\ \le\ \sum_{i>J}w_i+\sum_{i\in I_0}w_ir^k+\sum_{i\in I_1}w_i\,2\delta\ \le\ \varepsilon+r^k+\varepsilon\ \le\ 3\varepsilon
\]
once $r^k\le\varepsilon$. But $\sum_{b\in B_k}c\,\nulen(D_b)=c(1-2^{1-k})\ge c/2=4\varepsilon$ for $k\ge2$, a contradiction. ($\MF\ge2^{-O(1)}\nulen$ by \cref{thm:setting:universal}, so $\MF/\Mprod$ is unbounded.)
\end{proof}

For \cref{rem:versions:sharedsize}: the machine above certifies the labels of all inputs of length at most $d$ within $\omega_{1:d+1}$, so $\nulen^{(d+1)}=\nulen$ on such data, and \cref{prop:versions:tape}(a) with \cref{thm:versions:chain} gives $\Mprod(D)\ge2^{-O(1)}W(D)\ge2^{-\Kpad(d+1)-O(1)}\nulen(D)$.

\subsection{No coding theorem; the covering bound}\label{app:versions:nocoding}

\begin{proof}[Proof of \cref{thm:versions:nocoding}]
(i) Let $g:=\ell_d=\lceil\log(d/\log3)\rceil-1$, so that $3^{2^g}<2^d$ (c8(a) checks the arithmetic for all $d\le2\cdot10^5$ and along powers of $2$ up to $2^{200}$). Fewer than $2^g$ programs have length less than $g$; let $E$ be the set of their indices $V(p)$. Colour $x\in\{0,1\}^d$ by $(\chi_e(x))_{e\in E}\in\{0,1,*\}^E$ as in \cref{prop:versions:Wgap}. There are at most $3^{|E|}<2^d$ colours, so two distinct $x_1,x_2$ share a colour, and no $e\in E$ has $\varphi_e(x_1)=0$, $\varphi_e(x_2)=1$. So $\KF(D)\ge g\ge\log d-\log\log3-1>\log d-2$.
(ii) \Cref{prop:versions:tape}(b) with $|D|=2$.
(iii) The uniform product law, as in \cref{prop:versions:Wgap}.
For the last claim: by (i), $\log d\le\KF(D)+2$, so $\lceil\log(d+1)\rceil\le\KF(D)+4$; and $K(m)\le\log m+2\log\log m+O(1)$ for $m\ge2$.
\end{proof}

\begin{proof}[Proof of \cref{cor:versions:twopoint}]
Lower bound: \cref{thm:versions:nocoding}(i),(iii). Upper bound: $\LSFI(D)\ge0$, and \cref{prop:versions:covering} with $k=2$ gives $\KF(D)\le2+\log(2(d+2))+K(\lceil\log(d+2)\rceil,2)+O(1)$.
\end{proof}

\begin{proof}[Proof of \cref{prop:versions:covering}]
Put $s:=\lceil\log(d+2)\rceil$ and $d':=2^s-2\ge d$; then $d'+2<2(d+2)$ (c8(d)). There are at most $\binom{2^{d'+1}-1}{k}2^k\le2^{k(d'+2)}$ binary-labelled data sets of size $k$ on $\{0,1\}^{\le d'}$. Put $N_{\rm cov}:=\lfloor2^kk(d'+2)\ln2\rfloor+1$. $N_{\rm cov}$ independent uniform truth tables on $\{0,1\}^{\le d'}$ all miss a fixed such $D$ with probability $(1-2^{-k})^{N_{\rm cov}}<e^{-N_{\rm cov}2^{-k}}\le2^{-k(d'+2)}$. By the union bound the expected number of uncovered data sets is below $1$, so some family of $N_{\rm cov}$ truth tables covers all of them. The lexicographically first such family is computable from $(d',k)$, hence from $(s,k)$, by exhaustive search. A program consists of a $V$-program for $(s,k)$ followed by a $\lceil\log N_{\rm cov}\rceil$-bit index (its length is computable from $(s,k)$, so the program is self-delimiting); it outputs an index of the chosen truth table, undefined off $\{0,1\}^{\le d'}$. Its length is at most $K(s,k)+\log N_{\rm cov}+1+O(1)$, and $\log N_{\rm cov}\le k+\log(k(d'+2))\le k+\log(k(d+2))+1$. For the consequence, $\LSFI(D)\ge K(D\mid\dom D)-O(1)=k-\delta_D-O(1)$ by \cref{thm:setting:coding}(a).
\end{proof}

Check c8(c) computes, for $d\le3$ and $k\le3$, the exact expected number of uncovered data sets ($<1$ in all nine cases) and finds that random families of size $N_{\rm cov}$ cover everything in at least two thirds of the trials in every case. Explicit two-point coverings also exist: the $\lfloor N/2\rfloor$-subsets of $\{0,\dots,N-1\}$ that contain $0$, read as $N$ functions on $\binom{N-1}{\lfloor N/2\rfloor-1}$ inputs, show all four label patterns on every pair of inputs (c8(b), $N\le14$). This is the construction of \citet{kleitman1973families}, reported to be optimal; we did not check optimality against the source, and the checks confirm it only for $N\le6$ (c8(b)) and $N\le8$ (r11). Covering all inputs of length at most $d$ this way takes at most $d+6$ functions for $d\le40$ (c8(b)). Nothing above depends on optimality.

\subsection{Structured outputs}\label{app:versions:structured}

\begin{proof}[Proof of \cref{prop:versions:structured}]
A conditional prefix program for $y$ given $x$ is a monotone program, given $x$, that prints $y$; so $M(y\mid x)\ge2^{-c}\sum_{q:V(q,x)=y}2^{-|q|}\ge2^{-K(y\mid x)-c}$. A constant program prints $0^\infty$, so $M(0^d\mid\varepsilon)\ge2^{-c}$, while $K(0^d)=K(d)\pm O(1)$. For the end marker, ``$\ge$'' is the first claim applied to $y\$$, with $K(y\$\mid x)=K(y\mid x)\pm O(1)$. For ``$\le$'': the strings $y\$$, $y\in\{0,1\}^*$, are pairwise incomparable, so their cylinders are disjoint and $\sum_yM(y\$\mid x)\le1$; $y\mapsto M(y\$\mid x)$ is lsc uniformly in $x$, and the conditional coding theorem gives $M(y\$\mid x)\le2^{-K(y\mid x)+c}$.
\end{proof}
